% ===================================================================
%  TAULA N° 4 — Atge Madur e Vielhesa.
%  Traduction 2026 d'apres texto/io, controlee sur texto/fr et
%  texto/ca. Consigne : voir texto/oc/10-taula-01.tex.
% ===================================================================

%%K t04-apar-1 apar
\VUsaut{4.0mm}
\VUtitre{15.0pt}{60}{TAULA N° 4}
\VUsaut{1.6mm}
\VUtitre{11.4pt}{0}{Atge Madur e Vielhesa.}
\VUsaut{3.2mm}

%%K t04-tit-1 sub
\VUcentre{11.4pt}{0}{\textit{Primièra scèna.}}
\VUsaut{1.6mm}
\VUtitre{11.4pt}{0}{Lo Banquet de Nòça.}
\VUsaut{2.6mm}

%%K t04-tit-2 sub
\VUpk{10.2pt}{40}{L'auton.}
\VUsaut{3.0mm}

%%K t04-c1-01-1 p
1. --- Los joves an creissut. Un d'eles, Joan, se marida. Assistissèm al
\VUgras{banquet de nòça}, qu'amassa a l'entorn de la meteissa taula las familhas
dels dos espós.

%%K t04-c1-02-1 p
2. --- Sèm a l'\VUgras{auton}, al mes de \VUgras{setembre}. Lo vent freg
d'\VUgras{octòbre} e de \VUgras{novembre} a pas encara despolhat la campanha de
son vestit verd. L'aire del vèspre es tèbe e perfumat; per aquò lo sopar se
servís jos la \VUgras{galariá} \textsuperscript{(8)} que dona sul jardin.

%%K t04-c1-03-1 p
3. --- Al luènh, sul \VUgras{truc} \textsuperscript{(3)}, i a l'ostal dels parents
de Joan, un elegant \VUgras{castèl} \textsuperscript{(1)} escondut dins la verdura;
de bèlas vinhas, que donan un vin excellent, cobrisson lo pendís del puèg. Lo
vinhairon las cultiva e las suènha.

%%K t04-c1-04-1 p
4. --- Per dessús las vinhas passa una banda de \VUgras{perdises}
\textsuperscript{(2)}, espaurugadas per l'arribada del \VUgras{garda-caça}
\textsuperscript{(5)}. Aqueste, amb l'\VUgras{escopeta} jos lo braç e lo
\VUgras{carnièr} en bandolièra, bat los \VUgras{bartasses}
\textsuperscript{(4)} amb son \VUgras{gos} \textsuperscript{(6)}. Susvelha lo
domeni e empacha los braconièrs d'acabar amb la caça.

%%K t04-c1-05-1 p
5. --- Fòra de la \VUgras{galariá} s'enfila una \VUgras{trelha}, que ne pendolan
de \VUgras{rasims} \textsuperscript{(7)} pesucs.

%%K t04-c1-06-1 p
6. --- Las polidas flors d'auton qu'ornan la \VUgras{taula}
\textsuperscript{(16)} son de \VUgras{crisantèmas} \textsuperscript{(9)}; lo
\VUgras{paire} \textsuperscript{(10)} de la nòvia se n'es pausat una a la
botonièra del frac.

%%K t04-c1-07-1 p
7. --- Lo repais toca a sa fin. Lo \VUgras{sirvent} \textsuperscript{(11)}
comença de portar las postres e retirarà las diferentas pèças del cobèrt ---
\VUgras{culhièras} \textsuperscript{(14)}, \VUgras{forcetas}
\textsuperscript{(12)}, \VUgras{cotels} \textsuperscript{(13)},
\VUgras{plats} \textsuperscript{(15)} --- que servisson pas mai.

%%K t04-tit-3 sub
\VUpk{10.2pt}{40}{La Familha.}
\VUsaut{3.0mm}

%%K t04-c1-08-1 p
8. --- Tot lo mond sembla alègre, e lo sorire amb lo qual la \VUgras{maire}
\textsuperscript{(17)} del nòvi agacha sos enfants pròva son bonur.

%%K t04-c1-09-1 p
9. --- Aquesta nòça unís doas familhas riquas del país. Joan, lo \VUgras{marit}
\textsuperscript{(18)}, es \VUgras{filh} d'un grand proprietari e ven lo
\VUgras{gendre} d'un grand industrial.

%%K t04-c1-10-1 p
10. --- Los \VUgras{espós} son joves e çaquelà fasiá longtemps qu'èran
\VUgras{promeses}. La \VUgras{jova esposa} \textsuperscript{(19)}, Marta, es
encantadora amb son \VUgras{vestit blanc} ornat de flors d'irangièr.

%%K t04-c1-11-1 p
11. --- Son \VUgras{sògre} \textsuperscript{(20)} es fòrça content d'aver una
\VUgras{nòra} tan amabla, e la \VUgras{sògra} \textsuperscript{(21)} de Joan
sorís en veire sa \VUgras{filha} tan polida.

%%K t04-c1-12-1 p
12. --- A la cima de la taula s'assètan los autres \VUgras{convidats}
\textsuperscript{(22)}: \VUgras{parents, amics, cosins e cosinas}. Un \VUgras{cap
de servici} \textsuperscript{(23)}, amb la servieta jos lo braç, los servís amb
diligéncia.

%%K t04-tit-4 sub
\VUpk{10.2pt}{40}{Los Amics.}
\VUsaut{3.0mm}

%%K t04-c1-13-1 p
13. --- Al costat drech de la taula, vaicí una \VUgras{damaisèla}
\textsuperscript{(24)}, \VUgras{amiga} de la nòvia, e puèi lo \VUgras{fraire}
d'aquesta, qu'es son \VUgras{pairin de nòça} \textsuperscript{(25)}. Charra amb sa
\VUgras{dòna d'onor} \textsuperscript{(26)}, sòrre del nòvi, que per aquesta nòça
ven \VUgras{bèla-sòrre} de Marta. Es una jova encantadora. Pòrta de bèlas
\VUgras{jòias}, un \VUgras{colar de pèrlas} e un \VUgras{braçalet} d'aur.

%%K t04-c1-14-1 p
14. --- A prèp d'eles, lo sénher qu'es drech e auça sa copa es un \VUgras{amic}
\textsuperscript{(27)} de la familha. Es un dels \VUgras{testimònis del nòvi}.
Ofrís un \VUgras{brinde}, beu a la \VUgras{santat} dels nòvis e lor desira fòrça
bonur. Va en etiqueta, amb frac e \VUgras{sabatas de vernís}
\textsuperscript{(28)}. Acompanha la \VUgras{grand} \textsuperscript{(29)}, qu'es
\VUgras{veusa}.

%%K t04-c1-15-1 p
15. --- Lo jove de \VUgras{frac} \textsuperscript{(31)}, amb la mostacha prima e
negra, es lo \VUgras{bèu-fraire} \textsuperscript{(30)} de Joan. Es un enginhaire
distinguit. Es al costat d'una \VUgras{jova dòna} \textsuperscript{(33)} fòrça
eleganta, la \VUgras{sòrre màger} de Marta e maire de la polida dròlleta que ven,
al braç de son cosin, li ofrir una flor e li \VUgras{desirar} la \VUgras{bona
nuèch}. Aquesta dòna pòrta de \VUgras{pendents} fòrça bèles e un elegant
\VUgras{penchenatge}.

%%K t04-c1-16-1 p
16. --- Lo pichon cosin, tan graciós amb sa \VUgras{blòda}
\textsuperscript{(36)} e sas \VUgras{bragas} \textsuperscript{(37)} bofidas, fa
fòrça pietat. Es \VUgras{orfanèl} \textsuperscript{(35)}. Perdèt fòrça jove son
paire e sa maire. Son oncle es son \VUgras{tutor} \textsuperscript{(38)} e demorèt
solitari per elevar lo paure enfant. Es un òme excellent. Es un pauc calv e fòrça
miòp; pòrta de \VUgras{lunetas de pincèl}.

%%K t04-tit-5 sub
\VUsaut{2.4mm}
\VUpk{10.2pt}{40}{Las Postres.}
\VUsaut{3.0mm}

%%K t04-c1-17-1 p
17. --- Darrièr los convidats, sus una taula cobèrta d'una \VUgras{toalha}, son
preparadas las \VUgras{postres} \textsuperscript{(39)}. Dins una granda copa,
vaicí de fruchas: de \VUgras{persègas} \textsuperscript{(40)}, de \VUgras{peras}
\textsuperscript{(41)}, de \VUgras{poms} \textsuperscript{(42)}, d'\VUgras{abricòts}
\textsuperscript{(43)} e de \VUgras{rasim} \textsuperscript{(44)}. Al costat, de
\VUgras{ananàs} \textsuperscript{(45)}, un bèl \VUgras{gastèl}
\textsuperscript{(46)}, de \VUgras{botelhas de licor} \textsuperscript{(48)} e de
\VUgras{champanha} \textsuperscript{(49)}, e tanben de pilas de \VUgras{plats}
\textsuperscript{(47)} nets.

%%K t04-c1-18-1 p
18. --- Lo \VUgras{cavista} \textsuperscript{(50)} destapa una \VUgras{botelha}
\textsuperscript{(52)} de vin vièlh amb un \VUgras{tirabocon}
\textsuperscript{(51)}. Lo \VUgras{bocon} \textsuperscript{(53)} sembla fòrça
malaisit a tirar. Aqueste cavista pòrta una \VUgras{servieta}
\textsuperscript{(54)} jos lo braç.

%%K t04-c1-19-1 p
19. --- Aprèp lo sopar, los senhers passaràn al salon de fumar e las dònas
anaràn s'aprestar pel bal.

%%K t04-tit-6 sub
\VUsaut{5.0mm}
\VUcentre{11.4pt}{0}{\textit{Segonda scèna.}}
\VUsaut{1.6mm}
\VUpk{11.4pt}{40}{L'Anniversari del grand.}
\VUsaut{2.6mm}

%%K t04-tit-7 sub
\VUpk{10.2pt}{40}{La Visita.}
\VUsaut{3.0mm}

%%K t04-c2-01-1 p
1. --- Dètz ans son passats. Los parents de Joan an vielhit e aiman fòrça sos
tres felens, dos \VUgras{felens} \textsuperscript{(2)} e una \VUgras{felena}
\textsuperscript{(3)}. Coma uèi es l'\VUgras{anniversari del grand}, los enfants
venon lo felicitar.

%%K t04-c2-02-1 p
2. --- Lo pichon Pèire es encara fòrça mainat; a pas que tres ans. Vei que
s'apròcha lo \VUgras{cat} \textsuperscript{(1)}. Vòl careçar son bèl \VUgras{pel}
sedós e sa longa \VUgras{coa}; pensa pas que, jos aquelas graciosas
\VUgras{patas}, s'escondon de laidas onglas ponchudas que lo poirián plan
esgrafinhar.

%%K t04-c2-03-1 p
3. --- Sa sòrre Gabrièla, que li dona la man, es fòrça mai seriosa, e Marcèl, amb
son grand \VUgras{mantèl} \textsuperscript{(4)} e sa \VUgras{cencha}
\textsuperscript{(5)} de cuèr, sembla ja un omenet, malgrat las bragas cortas que
daissan veire sas \VUgras{caucetas} \textsuperscript{(6)}.

%%K t04-c2-04-1 p
4. --- Son paire s'es pausat son gròs \VUgras{pardessús} \textsuperscript{(7)}
d'ivèrn amb \VUgras{còl de pèl} \textsuperscript{(8)} e, dins la
\VUgras{pòcha} \textsuperscript{(9)}, pòrta un mocador. Sa femna pòrta son capèl
de feutre ornat de \VUgras{plumas} \textsuperscript{(10)}, sa \VUgras{vèsta}
\textsuperscript{(11)}, son \VUgras{boà de pèl} \textsuperscript{(12)} e son
\VUgras{manjon} \textsuperscript{(13)}.

%%K t04-c2-05-1 p
5. --- Fa fòrça freg, que l'ivèrn règna. La \VUgras{nèu} \textsuperscript{(19)} es
tombada tot lo jorn e los teulats dels ostals son blancs. Amb la \VUgras{nuèch},
lo freg ven encara mai viu. Al cèl, la \VUgras{luna} \textsuperscript{(17)} e las
\VUgras{estelas} \textsuperscript{(18)} lusisson e traucan l'\VUgras{escuresina}.

%%K t04-tit-8 sub
\VUsaut{4.2mm}
\VUpk{10.2pt}{40}{La Malautiá.}
\VUsaut{3.0mm}

%%K t04-c2-06-1 p
6. --- Lo paure \VUgras{bòrni} \textsuperscript{(20)} que traversa la plaça davant
la \VUgras{farmacia} \textsuperscript{(16)}, menat per son \VUgras{canich}
\textsuperscript{(21)}, es fòrça malastruc. Justina, la \VUgras{serviciala}
\textsuperscript{(14)}, pòrta de la cosina un \VUgras{bòl} \textsuperscript{(15)}
de \VUgras{tisana} fòrça cauda e plan sucrada.

%%K t04-c2-07-1 p
7. --- La \VUgras{grand} \textsuperscript{(23)}, que vòl anar quèrre sas
\VUgras{lunetas} \textsuperscript{(26)} sus la taula per legir l'\VUgras{ordonança}
\textsuperscript{(27)} que lo \VUgras{mètge} \textsuperscript{(25)} ven
d'escriure, se lèva de son fautuèlh en s'apiejant sus son \VUgras{baston}
\textsuperscript{(24)}.

%%K t04-c2-08-1 p
8. --- Sovent patís de pichonas \VUgras{indisposicions}, de \VUgras{vertolhs}, de
\VUgras{raumasses} e de \VUgras{mals de caissal}; a las cambas flacas e la vista
es pas mai gaire bona. Malgrat aquò, se trufa de bon grat de son vièlh
\VUgras{mètge}, que vòl totjorn li ordonar una \VUgras{pocion}
\textsuperscript{(28)}. Uèi es fòrça contenta d'aver sa jova familha a l'entorn
d'ela.

%%K t04-c2-09-1 p
9. --- Òm s'i tròba plan, dins la cambra plan barrada. Dins la \VUgras{chiminèa}
\textsuperscript{(29)} de mabre crèma un \VUgras{fuòc esplendid}
\textsuperscript{(30)}, e òm se sentís urós amb la doça \VUgras{lutz} de la
\VUgras{lampa} \textsuperscript{(31)} que venon d'alucar.

%%K t04-tit-9 sub
\VUsaut{4.2mm}
\VUpk{10.2pt}{40}{Lo Grand.}
\VUsaut{3.0mm}

%%K t04-c2-10-1 p
10. --- Quitament lo \VUgras{grand} \textsuperscript{(33)} oblida qu'es estat
malaut, que son \VUgras{reumatisme} lo ten clavat a son \VUgras{fautuèlh}
\textsuperscript{(34)} e que, ièr encara, aviá la \VUgras{fèbre e de
desfalhenças}. Se sovèn pas mai que l'ivèrn passat aguèt una \VUgras{pneumonia}
e qu'una \VUgras{tos} rauca e penosa lo fasiá fòrça patir.

%%K t04-c2-10-2 p
E çaquelà se remetèt amb fòrça pena. Ara va un pauc melhor. Mas la camba, que
pausa sus un \VUgras{coissin} \textsuperscript{(38)} pausat sus un pichon
\VUgras{tamboret} \textsuperscript{(39)}, li fa encara mal.

%%K t04-c2-11-1 p
11. --- Quand es plan embolhat dins sa cauda \VUgras{rauba de cambra}
\textsuperscript{(35)} amb de gròsses \VUgras{botons} \textsuperscript{(36)} de
nacre, e quand a a son costat, per li legir lo jornal, la pichona
\VUgras{orfanèla} \textsuperscript{(40)} amb sa \VUgras{còfa} blanca, sa
\VUgras{rauba} \textsuperscript{(42)} e sa \VUgras{pelerina}
\textsuperscript{(41)} blavas, se plang pas.

%%K t04-c2-12-1 p
12. --- Cada an, quand lo \VUgras{calendièr} \textsuperscript{(43)} torna portar
la \VUgras{data} del primièr de \VUgras{decembre}, espèra amb impaciéncia sos
\VUgras{felens}, que sap qu'oblidaràn pas aquela \VUgras{data} importanta.

%%K t04-c2-13-1 p
13. --- Se sovèn amb emocion dels temps luènhs ont el meteis anava felicitar lo
gloriós \VUgras{marescal} de l'Empèri, que son \VUgras{retrach}
\textsuperscript{(44)} pendola a la paret, e qu'es lo \VUgras{rèire-grand} dels
enfants.
\VUsaut{6.0mm}
\VUfilet{22mm}
